\documentclass[a4paper,10pt]{article}
\usepackage[margin=1.5cm]{geometry}
\usepackage{longtable}
\usepackage[T1]{fontenc}
\usepackage[utf8]{inputenc}

\title{Docker Cheatsheet}
\date{}

\begin{document}
\maketitle

\section*{1. Dockerfile Commands}

\begin{longtable}{|p{3cm}|p{11cm}|}
\hline
\textbf{Comando} & \textbf{Descripción} \\
\hline
FROM & Define la imagen base. Ej: \texttt{FROM ubuntu:22.04} \\
\hline
RUN & Ejecuta comandos durante el build. Ej: \texttt{RUN apt update \&\& apt install -y xclip} \\
\hline
COPY & Copia archivos del host al contenedor. Ej: \texttt{COPY . /app} \\
\hline
ADD & Similar a COPY pero permite URLs y extracción automática. \\
\hline
WORKDIR & Define directorio de trabajo. Ej: \texttt{WORKDIR /app} \\
\hline
CMD & Comando por defecto al iniciar el contenedor. \\
\hline
ENTRYPOINT & Ejecutable principal del contenedor. \\
\hline
ENV & Variables de entorno. Ej: \texttt{ENV PORT=8080} \\
\hline
EXPOSE & Documenta puerto usado. Ej: \texttt{EXPOSE 8000} \\
\hline
USER & Cambia el usuario activo. \\
\hline
VOLUME & Define punto de montaje. \\
\hline
ARG & Variable solo para build. \\
\hline
\end{longtable}

\section*{2. Comandos de Terminal}

\subsection*{Imágenes}

\begin{longtable}{|p{5cm}|p{9cm}|}
\hline
\texttt{docker build -t nombre .} & Construye una imagen \\
\hline
\texttt{docker images} & Lista imágenes locales \\
\hline
\texttt{docker rmi imagen} & Elimina imagen \\
\hline
\texttt{docker pull imagen} & Descarga imagen \\
\hline
\end{longtable}

\subsection*{Contenedores}

\begin{longtable}{|p{5cm}|p{9cm}|}
\hline
\texttt{docker run imagen} & Ejecuta contenedor \\
\hline
\texttt{docker run -it imagen bash} & Contenedor interactivo \\
\hline
\texttt{docker ps} & Lista contenedores activos \\
\hline
\texttt{docker ps -a} & Lista todos los contenedores \\
\hline
\texttt{docker stop id} & Detiene contenedor \\
\hline
\texttt{docker start id} & Inicia contenedor \\
\hline
\texttt{docker restart id} & Reinicia contenedor \\
\hline
\texttt{docker rm id} & Elimina contenedor \\
\hline
\texttt{docker logs id} & Muestra logs \\
\hline
\texttt{docker exec -it id bash} & Entra al contenedor \\
\hline
\texttt{docker inspect id} & Información detallada \\
\hline
\end{longtable}

\subsection*{Servicio Docker}

\begin{longtable}{|p{6cm}|p{8cm}|}
\hline
\texttt{sudo systemctl start docker} & Inicia servicio \\
\hline
\texttt{sudo systemctl stop docker} & Detiene servicio \\
\hline
\texttt{sudo systemctl restart docker} & Reinicia servicio \\
\hline
\texttt{sudo systemctl status docker} & Estado del servicio \\
\hline
\texttt{sudo systemctl enable docker} & Activa autoarranque \\
\hline
\texttt{sudo systemctl disable docker} & Desactiva autoarranque \\
\hline
\end{longtable}

\section*{3. Limpieza}

\begin{longtable}{|p{5cm}|p{9cm}|}
\hline
\texttt{docker system prune} & Limpia recursos no usados \\
\hline
\texttt{docker container prune} & Borra contenedores detenidos \\
\hline
\texttt{docker image prune} & Borra imágenes sin uso \\
\hline
\texttt{docker volume prune} & Borra volúmenes sin uso \\
\hline
\end{longtable}

\end{document}
